\chapter{Conclusion}
\label{ch:conclusion}

Automatic swara segmentation and classification play an important role in music information retrieval, digital music analysis, and computer-assisted music education. However, developing robust automatic transcription systems remains challenging due to the scarcity of accurately annotated continuous musical recordings. Manual annotation of note boundaries is a labor-intensive and time-consuming process, making it difficult to construct large datasets for training deep learning models. Motivated by this challenge, this research proposed an end-to-end framework that minimizes the dependence on manually annotated continuous recordings by utilizing synthetically generated melodies with automatically created ground-truth annotations.

The proposed framework consists of four major stages: dataset preparation, feature extraction, swara segmentation, and swara classification. An individual swara dataset containing recordings of the seven fundamental swaras was first collected from both publicly available recordings and self-recorded flute performances. These isolated recordings were then used to generate a synthetic melody dataset through randomized note selection, transition modeling, and multiple digital signal processing augmentations. Since the melody generation process automatically recorded the temporal boundaries of every inserted note, accurate frame-level annotations were obtained without requiring any manual labeling.

To evaluate the proposed framework, six different acoustic feature representations were investigated, namely Log-Mel, Pitch, RMS, Log-Mel + Pitch, Log-Mel + Pitch + RMS, and Constant-Q Transform (CQT). All experiments employed identical segmentation and classification architectures, ensuring that performance differences resulted solely from the selected feature representation. Experimental results demonstrated that the Constant-Q Transform provided the most effective representation for monophonic flute melodies, enabling the proposed framework to achieve the highest overall performance, with an end-to-end segmentation accuracy of 92.86\% and a note-level classification accuracy of 100\% on correctly segmented melodies.

Beyond the achieved performance, the primary contribution of this research lies in demonstrating that accurate swara segmentation and classification can be accomplished without relying on hundreds of hours of manually annotated continuous recordings. The proposed synthetic melody generation strategy significantly reduces the annotation effort while preserving reliable ground-truth information for model training and evaluation. This provides a practical and scalable alternative for developing automatic music transcription systems where annotated datasets are limited or unavailable.

Although the present work was validated only on flute recordings, the proposed methodology is not inherently restricted to a single instrument. Since the pipeline operates on isolated recordings of individual notes, the same synthetic data generation strategy can potentially be adapted to other monophonic musical instruments by replacing the source recordings while preserving the remaining framework. This flexibility provides a promising foundation for future research in automatic music transcription and music information retrieval.

In conclusion, this thesis demonstrates that combining synthetic dataset generation with deep learning-based segmentation and classification models offers an effective solution for automatic swara identification. The proposed framework not only achieves promising performance but also presents a practical approach for reducing manual annotation requirements, making it a valuable contribution toward scalable and data-efficient research in computational music analysis.
